\chapter{Introduction}
\label{ch:introduction}

The healthcare industry today faces significant challenges in managing and accessing patient medical information. In the current system, patient data is often scattered across multiple hospitals and healthcare providers, creating fragmented medical histories that are difficult to access when needed. This fragmentation leads to several critical issues: patients become dependent on physical reports, undergo repeated diagnostic tests, face limited options for remote consultations, and encounter difficulties in obtaining certain medications, especially rare ones. Additionally, hospitals operate without a common digital platform to efficiently manage patient data and streamline interactions between healthcare providers and patients.

\noindent
The digital transformation of healthcare has become increasingly important in recent years, with the COVID-19 pandemic highlighting the urgent need for remote healthcare solutions. Patients, particularly the elderly and those with mobility issues, require accessible healthcare services that do not always necessitate physical hospital visits. Furthermore, the need for a centralized, secure system to store and access medical records has become paramount in ensuring continuity of care and improving healthcare outcomes.

\section{Problem Statement}
The existing healthcare ecosystem suffers from the following key problems:
\begin{itemize}
    \item \textbf{Fragmented Medical Records}: Patient medical information is distributed across different hospitals, making it difficult to obtain complete medical histories when required for diagnosis or treatment.
    \item \textbf{Dependency on Physical Reports}: Patients must carry physical copies of their medical reports, which can be lost, damaged, or forgotten during emergencies.
    \item \textbf{Repeated Diagnostic Tests}: Due to inaccessible medical histories, patients often undergo redundant tests, increasing healthcare costs and inconvenience.
    \item \textbf{Limited Remote Consultation Options}: Many patients, especially the elderly and those in remote areas, struggle to access healthcare services that require physical presence.
    \item \textbf{Difficulty in Obtaining Medications}: Patients face challenges in locating and purchasing rare or specific medicines that may only be available at certain hospital pharmacies.
    \item \textbf{Lack of a Common Digital Platform}: Hospitals lack a unified digital interface to efficiently manage patient data, appointments, consultations, and billing processes.
\end{itemize}

\section{Objectives}
The primary objectives of the BondHealth project are:

\begin{enumerate}
    \item \textbf{Develop a Centralized Digital Platform}: Create a web-based healthcare platform that connects hospitals and patients within a specific geographical area, functioning as a digital interface for healthcare services.
    
    \item \textbf{Implement Digital Patient Identities}: Provide unique, secure digital IDs to patients (similar to DigiLocker) for storing and accessing medical records including consultation history, diagnostic reports, and medical scans.
    
    \item \textbf{Enable Hospital Discovery}: Display registered hospitals along with available departments and doctor information, allowing patients to make informed decisions about their healthcare providers.
    
    \item \textbf{Facilitate Remote Consultations}: Support scheduled chat and video call consultations between patients and doctors, making healthcare accessible to those unable to visit hospitals in person.
    
    \item \textbf{Integrate Pharmacy Services}: Enable online purchase of medicines, including rare medications, from hospital pharmacies with electronic billing for transparent, paperless transactions.
    
    \item \textbf{Ensure Security and Privacy}: Implement robust security measures to protect sensitive patient data and ensure compliance with healthcare data protection standards.
\end{enumerate}

\section{Scope of the Project}
The scope of BondHealth encompasses the following aspects:

\subsection{Geographical Scope}
The platform initially targets a specific geographical area, connecting hospitals and patients within that region. This focused approach allows for better management, testing, and refinement of the system before potential expansion to wider areas.

\subsection{Functional Scope}
\begin{itemize}
    \item \textbf{Patient Registration and Authentication}: Secure registration process assigning unique patient IDs with robust authentication mechanisms.
    \item \textbf{Medical Records Management}: Digital storage and retrieval of consultation histories, diagnostic reports, and medical scans.
    \item \textbf{Hospital and Doctor Directory}: Comprehensive listing of registered hospitals with department-wise doctor information.
    \item \textbf{Teleconsultation Services}: Scheduling and conducting remote consultations through chat and video calls.
    \item \textbf{Online Pharmacy}: E-commerce functionality for medicine purchases from hospital pharmacies.
    \item \textbf{Electronic Billing}: Digital invoice generation and payment processing for consultations and medicine purchases.
\end{itemize}

\subsection{Technical Scope}
The project involves developing a responsive web application using modern web technologies, implementing secure database management for medical records, integrating video calling APIs for teleconsultation, and ensuring data encryption and privacy compliance.
